\scanpage{315}
\mspage{14}
% archive sheet number: 316

dans $\widehat\Delta_{\mathrm{ab}}=\mathit{Ab}_{\ast}$ est \uline{universel} pour les
complexes de cochaînes dans des cat.
additives $A$ avec $\varinjlim$ qc. (\uline{NB} Par DP
on peut aussi interpréter $\mathit{Ab}_{\ast}$ comme $\mathit{Ab}_{\cdot}$
dans cette description de la cat. contenant
l'objet universel --).

b) En dualisant et interprétant
$n\longmapsto\mathbb C_{\ast}(\Delta^{n}_{\ast})$ comme objet \uline{simplicial}
dans $(\mathit{Ab}_{\ast})^{\circ}$\inserted{$=(\mathit{Ab}^{\circ})^{\ast}$}, regardée comme cat additive
avec $\varprojlim$, il est universel pour les
cat. $B$ avec $\varprojlim$ et foncteurs entre
icelles y commutant. De $=$ pour le
complexe des chaînes associé.

Du point de vue de l'Algèbre universelle,
\deleted{cela veut dire que les constructions des} \inserted{on peut exprimer le contenu de} ces
réflexions de la façon suivante

1. Les inv. contravariants ``commutant aux
$\varprojlim$'' en $X_{\ast}\in\mathit{Ss}$, à valeurs dans une
cat. \uline{additive} $A$ avec $\varprojlim$, s'expriment
\unsure{au choix|en deux} en termes d'objets simpliciaux
de $A$, ou de complexes de chaînes de $A$.

2. Les constructions qu'on peut faire sur de
tels invariants (\inserted{pour $A$ variable}, de façon compatible aux
foncteurs $A\longrightarrow A'$ commutant aux $\varprojlim$)
s'expriment à l'aide des objets
de $(\mathit{Ab}_{\ast})^{\circ}$ (où se trouve l'objet
simplicial -- pour cat. add, $\varprojlim$ --
\uline{universel}), ou de $(\mathit{Ab}_{\cdot})$

et 1′), 2′) les énoncés duals (formulés
à partir de 1), 2))

Si $F,K_{\ast},K_{\cdot}$ se correspondent dans
1, on a
\[
\left\{\begin{aligned}
F(X_{\ast})&\simeq\operatorname{Hom}_{\ast}(\mathbb C_{\ast}(X_{\ast}),K_{\ast})
\simeq\operatorname{Hom}_{\cdot}(\mathbb C_{\cdot}(X_{\ast}),K_{\cdot})\\
K_{\cdot}&=\mathrm{ND}_{\cdot}(K_{\ast}),\quad K_{\ast}=\mathrm{DP}_{\ast}(K_{\cdot})
\end{aligned}\right.
\]

Si de $\overline{\mathrm m}$ $\Omega$ est une ``opération sur des
invariants'', et si $L_{\ast}$\inserted{$\in\mathit{Ab}_{\ast}$} et $L_{\cdot}\in\mathit{Ab}_{\cdot}$ y
sont associés, on aura
\margin{$\Omega$ a deux \unsure{aspects|expressions} ``opérations sur des $K_{\ast}$'' ($\Omega^{\ast}$) et ``op. sur des $K_{\cdot}$'' ($\Omega^{\cdot}$) -- pour trouver des objets de $B$ \deleted{\unsure{provenant}} à partir d'un complexe (simp. ou de chaînes) dans $B$}
\[
\left\{\begin{aligned}
\Omega(F)&\simeq\operatorname{Hom}_{\ast}(L_{\ast},K_{\ast})\inserted{=\Omega_{\ast}(K_{\ast})}
\simeq\operatorname{Hom}_{\cdot}(L_{\cdot},K_{\cdot})\inserted{=\Omega_{\cdot}(K_{\cdot})}\\
L_{\cdot}&=\mathrm{ND}_{\ast}(L_{\ast}),\quad L_{\ast}=\mathrm{DP}_{\ast}(L_{\cdot})
\end{aligned}\right.
\]

(La signification des symboles $\operatorname{Hom}_{\ast}$ (resp $\operatorname{Hom}_{\cdot}$)
pour \deleted{\unsure{t}} un $C_{\ast}$ ou $L_{\ast}$ \inserted{dans $\mathit{Ab}_{\ast}$} (resp. $C_{\cdot}$ ou $L_{\cdot}$\inserted{dans $\mathit{Ab}_{\cdot}$})
a été expliquée plus haut).

Dualement, \inserted{pour} les \deleted{\unsure{opérations}} \inserted{invariants} \inserted{d'objets $X_{\ast}\in\mathit{Ss}$}, covariants
commutant aux $\varinjlim$, à valeurs dans
$A$, si $G,K^{\ast},K^{\cdot}$ se correspondent,
on aura
\margin{à préciser la signification de ces formules}
\[
\left\{\begin{aligned}
G(X_{\ast})&\simeq\deleted{\unsure{Z^{0}_{\ast}}}(\mathbb C_{\ast}(X_{\ast})\overset{\ast}{\otimes}K_{\deleted{\unsure{\ast}}}^{\ast})
\simeq\deleted{\unsure{Z^{0}_{\ast}}}(\mathbb C_{\cdot}(X_{\ast})\overset{\cdot}{\otimes}K^{\cdot})\\
K_{\unsure{\cdot}}^{\cdot}&=\mathrm{ND}^{\cdot}(K^{\ast}),\quad K^{\ast}=\mathrm{ND}^{\ast}(K^{\cdot})
\end{aligned}\right.
\]

et si l'opération $\Omega$, $L_{\ast}$ et $L_{\cdot}$ se
correspondent, ($L_{\ast}\in\mathit{Ab}_{\ast}$, $L_{\cdot}\in\mathit{Ab}_{\cdot}$)
on aura
\[
\left\{\begin{aligned}
\Omega(G)&\simeq L_{\ast}\overset{\ast}{\otimes}K^{\ast}
\simeq L_{\cdot}\overset{\cdot}{\otimes}K^{\cdot}\\
L_{\cdot}&\simeq\mathrm{ND}_{\cdot}(L_{\ast}),\quad L_{\ast}\simeq\mathrm{DP}_{\ast}(L_{\cdot})
\end{aligned}\right.
\]
